\section{Related Work}
\label{sec:related}

\para{Generative monoculture.}
\citet{wu2024monoculture} formalize generative monoculture as a distribution shift in which LLM outputs are statistically narrower than the human-generated source data: $\text{Dispersion}(h(x) | x \sim P_{\text{gen}}) < \text{Dispersion}(h(x) | x \sim P_{\text{src}})$. They demonstrate this across book reviews (sentiment clustering near positive) and code solutions (reduced algorithm variety), showing that changing temperature, sampling strategies, and prompts is insufficient to overcome the effect. \citet{correlated2025errors} extend this finding, showing that errors across different LLMs are more correlated with each other than with ground truth, and that monoculture worsens within model families and with RLHF training. Unlike these works, which focus on objective or semi-objective tasks, we test monoculture in purely subjective evaluative judgment---where there is no single correct answer, making convergence especially revealing.

\para{Creative homogeneity across LLMs.}
\citet{wenger2025homogeneity} apply standardized creativity tests---the Alternative Uses Task, Forward Flow, and Divergent Association Task---to multiple LLMs and humans. They find that while individual LLM responses score competitively on creativity metrics, the population-level semantic similarity across LLM responses far exceeds that of human responses. Even creativity-oriented system prompts produce only marginal diversity gains. Our work complements this finding by testing a qualitatively different task: rather than generating creative uses for common objects, we ask models to exercise \textit{taste}---to identify items that are undervalued relative to their quality, which requires reasoning about both quality and popularity simultaneously.

\para{LLM creativity benchmarks.}
Several benchmarks evaluate LLM creativity along different dimensions. \citet{divergent2024creativity} benchmark LLMs against 100K+ human responses on the Divergent Association Task and creative writing, finding that LLMs match average human performance but fall short of highly creative individuals. \citet{hou2026creativityprism} propose CreativityPrism, decomposing creativity into quality, novelty, and diversity across 8 tasks with 17 metrics, and find that novelty shows weak or negative correlations with other creativity dimensions---high novelty in one domain does not transfer. \citet{liveideabench2024} introduce LiveIdeaBench with minimal-context prompts for divergent thinking evaluation. These benchmarks use abstract stimuli (word association, object uses) rather than domain-specific knowledge, leaving open whether LLMs can generate novel judgments in knowledge-rich domains.

\para{Alignment and the novelty--appropriateness tradeoff.}
\citet{beyond2026divergent} introduce the Conditional Divergent Association Task and find that alignment training shifts models toward ``appropriate'' responses at the cost of novelty---smaller model families sometimes show more creativity than larger, more aligned ones. \citet{ouyang2022training} describe how RLHF training rewards responses matching broad human preferences, which naturally pushes outputs toward consensus. This tradeoff directly explains why LLMs produce ``safe'' underrated picks: RLHF training optimizes for responses that most annotators would approve, which are by definition the popular choices.

\para{Content homogenization.}
Beyond creativity-specific work, \citet{padmakumar2024diversity} show that co-writing with LLMs reduces content diversity across users, and \citet{anderson2024homogenization} demonstrate homogenization effects in creative ideation when humans use LLM assistance. Our diagnostic extends this line of work from human--LLM collaboration to LLM-only generation, measuring how models perform when asked to exercise independent judgment without human guidance.
